\section{问题重述}
题目给出一条由223节板凳组成的“板凳龙”，要求在盘入、调头和盘出三个阶段中，结合把手约束、板凳宽度与调头空间边界，完成位置、速度、终止时刻、最小螺距、调头曲线和最大安全速度等计算。

从建模角度看，问题的核心不在于单个局部几何量，而在于如何把“龙头沿路径运动”与“整队板凳的相邻把手约束”统一到同一套参数化框架中。为此，本文先以等距螺线刻画盘入路径，再用逐节递推恢复全队状态；随后把每节板凳视作有限宽度矩形，利用碰撞检测判定运动可行性；最后围绕调头曲线长度与速度上限，分别转化为两圆弧几何构造和速度缩放问题。

\textbf{问题一：} 在螺距为0.55 m、龙头前把手速度恒为1 m/s的条件下，给出0--300 s内全部把手的位置和速度。

\textbf{问题二：} 沿问题一的盘入螺线继续运动，求首次发生碰撞前的终止时刻，并输出终止状态。

\textbf{问题三：} 在半径为4.5 m的调头空间内，求使龙头前把手能够盘入边界的最小螺距。

\textbf{问题四：} 在螺距为1.7 m的条件下，以螺线与调头圆的交点为调头入点，构造调头空间内的S形连接曲线。

\textbf{问题五：} 在问题四路径上，求满足所有把手速度均不超过2 m/s的龙头最大恒定速度。

\section{问题分析}
五个问题共享同一条主线：先确定龙头轨迹，再由把手间的几何约束回推整队状态，最后用碰撞与速度条件筛选可行解。

对问题一和问题二，关键在于建立“弧长--参数--坐标”的转换。龙头沿等距螺线运动时，弧长是时间的线性函数；只要能解出当前时刻对应的螺线参数，就能顺次恢复每节板凳的把手位置。问题二在此基础上增加了板凳实体宽度，因此不能只看把手中心线，而必须将每节板凳近似为有厚度的矩形，并检测非相邻板凳之间是否发生接触。

对问题三，最小螺距问题本质上是一个单变量可行性搜索。螺距越小，螺线越紧，整队越容易在调头空间边界附近产生干涉；因此可以把“给定螺距是否可行”作为判定函数，再用二分搜索逼近最小可行值。

对问题四，调头曲线由两段相切圆弧组成，且半径满足比例约束。入点由螺线与调头圆交点确定后，切线方向、圆弧半径和圆心角均可由几何关系推出。

对问题五，路径几何固定后，头部恒速会线性传递到全队各把手的速度幅值。于是只需先计算单位头速下的速度放大倍数，再通过“最大值不超过2 m/s”反推出头速上界即可。
